% =====================================================================
\section{Introduction}
% =====================================================================

\paragraph{Motivation and context.}
After fabrication in the clean room, devices have to be tested electrically to check
whether the process worked as expected. For this, simple test structures are used, such as
MIS (metal--insulator--semiconductor) capacitors and p--n junctions. They contain the main
elements of a MOS transistor (gate oxide, Si/SiO$_2$ interface, doped substrate, junctions)
but are simple enough to be analysed with basic semiconductor physics.
Capacitance--voltage (C--V) measurements on an MIS capacitor give the oxide thickness
$t_{ox}$, the substrate doping $N_A$, the flat-band voltage $V_{FB}$ and the charge density
$N_s$ at the Si/SiO$_2$ interface. Current--voltage (I--V) measurements on diodes give the
ideality factor, the series resistance, the leakage current and the breakdown voltage.
By comparing these values with the targets of the process, we can say which fabrication steps
went well and which ones should be improved.

\paragraph{Devices.}
The structures come from the MOSTEC layout made in the CIME Nanotech clean room: a p-type
Si substrate, a thermal oxide (target thickness \SIrange{45}{50}{nm}) and an n$^+$ poly-Si
gate for the MIS capacitors. The n$^+$/p diodes exist in two shapes with almost the same area
($\approx\SI{4e-4}{cm^2}$): a square one (small perimeter) and a rake-shaped one (large
perimeter). In both cases the top contact is taken with a probe tip and the p substrate
through the chuck.

\paragraph{Experimental means.}
The measurements are done on a probe station in a dark box. C--V curves are measured with an
Agilent E4980A LCR meter controlled by the MDC software, which applies a DC voltage ramp with a
small AC signal and records the capacitance $C$ and the conductance $G$, at \SI{1}{kHz} or
\SI{1}{MHz}. I--V curves are measured with a semiconductor parameter analyser (SMU).
All measurements were done at room temperature (\SI{23}{\celsius}).

\paragraph{Organization of the report.}
\Cref{sec:diodes} deals with the diodes, \cref{sec:mis} with the MIS capacitor and
\cref{sec:conclusion} summarizes the results.
Constants used: $\varepsilon_0=\SI{8.85e-14}{F/cm}$, $\varepsilon_{ox}=3.9\,\varepsilon_0$,
$\varepsilon_{Si}=11.7\,\varepsilon_0$, $k_BT/q=\SI{26}{mV}$,
$n_i=\SI{1.45e10}{cm^{-3}}$.
